% GENERATED FILE. Edit canonical lessons, then run npm run generate:books.
% canonical-source-hash: fnv1a64:0bd5dfc6ccacebc5
% canonical-lessons: RU-C461-polyana, RU-C461-opushka, RU-C461-chashcha, RU-C461-mokh, RU-C461-shishka

\chapter{In the forest}
\label{ch:ru-in-the-forest}
\hlchaptermodality{461}

\begin{chapteropening}
\textbf{By the end of this chapter:} \emph{I can describe a forest: a clearing, the edge of the forest, a thicket, moss and a cone.}

Complete the last lesson of chapter 461: I can describe a forest: a clearing, the edge of the forest, a thicket, moss and a cone.
\end{chapteropening}

\section[\texorpdfstring{\ru{поляна} (polyána)}{поляна (polyána)}]{\ru{поляна} (polyána) — a clearing, a glade}

\label{lesson:RU-C461-polyana}



\begin{quote}
Before the new one: say the Russian for a pond, then the Russian for to babble.
\end{quote}

\subsection*{You'll want to know: \ru{поляна}}
\textbf{\ru{поляна}} (\emph{polyána}) — \textquotedblleft{}a clearing, a glade\textquotedblright{}.

\begin{quote}
\ru{На} \ru{поляне} \ru{мы} \ru{увидели} \ru{оленя}.
\end{quote}

\begin{quote}
\emph{Na polyáne my uvídeli olénya.}
\end{quote}

\begin{quote}
We saw a deer in the clearing.
\end{quote}

Feminine. Built on \emph{\ru{поле}} — a field.

\begin{cousinweb}
\emph{\ru{поляна}} is built on \emph{\ru{поле}} 'a field, open land' with \emph{-\ru{яна}}. EtymDB lists \emph{\ru{поляна}} without a deeper source.
\end{cousinweb}

\begin{grammarlens}[title={What you've built}]
One more word for telling, explaining and describing at B1.
\end{grammarlens}

\subsection*{Your turn}
\begin{itemize}
\raggedright
  \item \emph{Say it:} \emph{polyána}
  \item \emph{Say it:} \emph{polyána}, once more
  \item \emph{Say it:} a sentence of your own with \emph{polyána}
  \item \emph{Recall:} say the Russian for a pond, then the Russian for to babble, then say \emph{polyána} again
\end{itemize}

\begin{tcolorbox}[breakable,colback=teal!4,colframe=teal!35!black,title={Before you move on}]
What does \ru{поляна} mean? (\textquotedblleft{}A clearing, a glade\textquotedblright{}.) Say it once more.
\end{tcolorbox}
\section[\texorpdfstring{\ru{опушка} (opúshka)}{опушка (opúshka)}]{\ru{опушка} (opúshka) — the edge of a forest}

\label{lesson:RU-C461-opushka}



\begin{quote}
Before the new one: say the Russian for to babble, then the Russian for a clearing.
\end{quote}

\subsection*{You'll want to know: \ru{опушка}}
\textbf{\ru{опушка}} (\emph{opúshka}) — \textquotedblleft{}the edge of a forest\textquotedblright{}.

\begin{quote}
\ru{На} \ru{опушке} \ru{леса} \ru{стоит} \ru{маленький} \ru{дом}.
\end{quote}

\begin{quote}
\emph{Na opúshke lésa stoít málenkiy dom.}
\end{quote}

\begin{quote}
There is a small house at the edge of the forest.
\end{quote}

Feminine. \emph{\ru{На} \ru{опушке}} — at the forest edge.

\begin{cousinweb}
\emph{\ru{опушка}} is usually linked with \emph{\ru{опушить}} 'to edge with fur', on \emph{\ru{пух}} 'down': the forest's fringe. EtymDB lists the word without a source, so the line is hedged.
\end{cousinweb}

\begin{grammarlens}[title={What you've built}]
One more word for telling, explaining and describing at B1.
\end{grammarlens}

\subsection*{Your turn}
\begin{itemize}
\raggedright
  \item \emph{Say it:} \emph{opúshka}
  \item \emph{Say it:} \emph{opúshka}, once more
  \item \emph{Say it:} a sentence of your own with \emph{opúshka}
  \item \emph{Recall:} say the Russian for to babble, then the Russian for a clearing, then say \emph{opúshka} again
\end{itemize}

\begin{tcolorbox}[breakable,colback=teal!4,colframe=teal!35!black,title={Before you move on}]
What does \ru{опушка} mean? (\textquotedblleft{}The edge of a forest\textquotedblright{}.) Say it once more.
\end{tcolorbox}
\section[\texorpdfstring{\ru{чаща} (cháshcha)}{чаща (cháshcha)}]{\ru{чаща} (cháshcha) — a thicket, the depths of a forest}

\label{lesson:RU-C461-chashcha}



\begin{quote}
Before the new one: say the Russian for a clearing, then the Russian for the edge of a forest.
\end{quote}

\subsection*{You'll want to know: \ru{чаща}}
\textbf{\ru{чаща}} (\emph{cháshcha}) — \textquotedblleft{}a thicket, the depths of a forest\textquotedblright{}.

\begin{quote}
\ru{Медведь} \ru{ушёл} \ru{в} \ru{чащу}.
\end{quote}

\begin{quote}
\emph{Medvéd ushyól v cháshchu.}
\end{quote}

\begin{quote}
The bear went off into the thicket.
\end{quote}

Feminine. \emph{\ru{В} \ru{лесной} \ru{чаще}} — deep in the forest.

\begin{cousinweb}
\emph{\ru{чаща}} is inherited from Proto-Slavic \emph{čęsťa} (EtymDB), built on the root of \emph{\ru{частый}} 'dense, frequent': where the trees stand thick.
\end{cousinweb}

\begin{grammarlens}[title={What you've built}]
One more word for telling, explaining and describing at B1.
\end{grammarlens}

\subsection*{Your turn}
\begin{itemize}
\raggedright
  \item \emph{Say it:} \emph{cháshcha}
  \item \emph{Say it:} \emph{cháshcha}, once more
  \item \emph{Say it:} a sentence of your own with \emph{cháshcha}
  \item \emph{Recall:} say the Russian for a clearing, then the Russian for the edge of a forest, then say \emph{cháshcha} again
\end{itemize}

\begin{tcolorbox}[breakable,colback=teal!4,colframe=teal!35!black,title={Before you move on}]
What does \ru{чаща} mean? (\textquotedblleft{}A thicket, the depths of a forest\textquotedblright{}.) Say it once more.
\end{tcolorbox}
\section[\texorpdfstring{\ru{мох} (mokh)}{мох (mokh)}]{\ru{мох} (mokh) — moss}

\label{lesson:RU-C461-mokh}



\begin{quote}
Before the new one: say the Russian for the edge of a forest, then the Russian for a thicket.
\end{quote}

\subsection*{You'll want to know: \ru{мох}}
\textbf{\ru{мох}} (\emph{mokh}) — \textquotedblleft{}moss\textquotedblright{}.

\begin{quote}
\ru{Камни} \ru{у} \ru{ручья} \ru{покрыты} \ru{мхом}.
\end{quote}

\begin{quote}
\emph{Kámni u ruchyá pokrýty mkhom.}
\end{quote}

\begin{quote}
The stones by the stream are covered in moss.
\end{quote}

Masculine; the \ru{о} drops out: \emph{\ru{мха}}, \emph{\ru{во} \ru{мху}}.

\begin{cousinweb}
\emph{\ru{мох}} is inherited from Proto-Slavic \emph{mŭxŭ} 'moss' (EtymDB), which EtymDB takes back to a PIE root \emph{mews-}.
\end{cousinweb}

\begin{grammarlens}[title={What you've built}]
One more word for telling, explaining and describing at B1.
\end{grammarlens}

\subsection*{Your turn}
\begin{itemize}
\raggedright
  \item \emph{Say it:} \emph{mokh}
  \item \emph{Say it:} \emph{mokh}, once more
  \item \emph{Say it:} a sentence of your own with \emph{mokh}
  \item \emph{Recall:} say the Russian for the edge of a forest, then the Russian for a thicket, then say \emph{mokh} again
\end{itemize}

\begin{tcolorbox}[breakable,colback=teal!4,colframe=teal!35!black,title={Before you move on}]
What does \ru{мох} mean? (\textquotedblleft{}Moss\textquotedblright{}.) Say it once more.
\end{tcolorbox}
\section[\texorpdfstring{\ru{шишка} (shíshka)}{шишка (shíshka)}]{\ru{шишка} (shíshka) — a cone (of a pine or a fir); a bump}

\label{lesson:RU-C461-shishka}



\begin{quote}
Before the new one: say the Russian for a thicket, then the Russian for moss.
\end{quote}

\subsection*{You'll want to know: \ru{шишка}}
\textbf{\ru{шишка}} (\emph{shíshka}) — \textquotedblleft{}a cone (of a pine or a fir); a bump\textquotedblright{}.

\begin{quote}
\ru{Под} \ru{елью} \ru{лежат} \ru{шишки}.
\end{quote}

\begin{quote}
\emph{Pod yélyu lezhát shíshki.}
\end{quote}

\begin{quote}
There are cones lying under the fir tree.
\end{quote}

Feminine. Also a bump on the head; colloquial \emph{\ru{большая} \ru{шишка}} — a big shot.

\begin{cousinweb}
\emph{\ru{шишка}} is inherited through Old East Slavic \emph{šišĭka} from Proto-Slavic \emph{šišĭka} (EtymDB).
\end{cousinweb}

\begin{grammarlens}[title={What you've built}]
One more word for telling, explaining and describing at B1.
\end{grammarlens}

\subsection*{Your turn}
\begin{itemize}
\raggedright
  \item \emph{Say it:} \emph{shíshka}
  \item \emph{Say it:} \emph{shíshka}, once more
  \item \emph{Say it:} a sentence of your own with \emph{shíshka}
  \item \emph{Recall:} say the Russian for a thicket, then the Russian for moss, then say \emph{shíshka} again
\end{itemize}

\begin{tcolorbox}[breakable,colback=teal!4,colframe=teal!35!black,title={Before you move on}]
What does \ru{шишка} mean? (\textquotedblleft{}A cone (of a pine or a fir); a bump\textquotedblright{}.) Say it once more.
\end{tcolorbox}
